\scanpage{90}
\mspage{77}

d'éléments de $C'$, donc $\unsure{u|v}^*(C')\subset C$ implique $u^*(\widetilde C')\subset C$ puisque $C$ stable par fact. dir. Prouvons enfin $u^*(\mathrm{T}\widetilde C')\subset\mathrm{TC}$\inserted{($\overset{\mathrm{déf}}{=}C\cap W$)}. \deleted{\unsure{Comme}} $\mathrm{T}\widetilde C'=\widetilde C'\cap W'$, cela résulte de $u^*(\widetilde C')\subset C$, et de $u^*(W')\subset W$.

Prouvons b$'$). La stabilité de $C'$\inserted{ et de $\widetilde C'$,} par composition transfinie est acquise, il reste à prouver \deleted{\unsure{que}} stabilité de $\mathrm{TC}'=C'\cap W'$ et de $\mathrm{T}\widetilde C'=\widetilde C'\cap W'$. Il suffit de voir que le composé transfini \inserted{$f$} en question est dans $W'$, donc que $u^*(f)\in W$, mais cela résulte du fait que $u^*$ commute à la compos. transfinie par (4.5.4), et que $u^*$ transforme $\mathrm{TC}'$ et $\mathrm{T}\widetilde C'$ dans $\mathrm{TC}$.

\uline{Prouvons c}\inserted{ pour $C'$ (donc aussi pour $\widetilde C'$)}. Soit $C_0$ une petite partie de $C$ telle que $C\subset\widetilde C_0$, i.e $C$ est contenu dans la clôture de $C_0$ par les trois op. de \unsure{coch.} de base, de composition transfinie, de fact. directs. On voit alors que $u_!(C)$ est dans la clôture de $u_!(C_0)$ pour \unsure{les|ces} opérations, donc aussi $C'$, et on peut donc prendre $C'_0=u_!(C_0)$.

\uline{Prouvons enfin d)}, c'est la $=$ assertion pour $C'$ et pour $\widetilde C'$, puisque
\[
C'_*= \widetilde C'_* \bigl(\overset{\mathrm{déf}}{=}\mathrm{TF}'\bigr),
\]
% archive sheet number: 8[3 ou 9]
